% *- coding: UTF-8 -*-
% !TEX builder = Texlive
% !TEX prggram = xelatex

%文档设置及格式
\documentclass[UTF8,a4paper,12pt,titlepage]{ctexart}
\include{document_setting}

% 设置文件属性内容
\hypersetup{pdfauthor={吴毓明},
            pdftitle={次氯酸钠加药装置操作说明},
            pdfsubject={使用说明},
            pdfkeywords={任何问题，联系61508712@qq.com}}

% 定义页眉页脚
\include{foot_head}

% 表格格式定义
\include{table_format}

% 图片
\include{picture_configuration}

% 特殊符号(例如千分号等)
\usepackage{textcomp}

% 流程图
\include{flow_chart_preamble}

% 开始正文
\begin{document}

% 制作封面
\title{
   \zihao{1} 次氯酸钠加药装置（单投加点）\\[2.5mm]
   \zihao{1} 操作说明\\
   {\noindent}\rule{16cm}{1pt}\\[3mm]
   \zihao{2} 宁夏智诚博创自动化科技有限公司\\
   % \zihao{4} 吴毓明\\
   %\zihao{4} Wym\\
   %\zihao{4} 2022年12月27日
   \begin{figure}[h]
      \centering
      \includegraphics[height=10cm]{g11.PNG}
   \end{figure}
}

\maketitle % 封面



\tableofcontents % 目录

\newpage % 换页

%\part{系统组成}
%ctex的\chapter{sdsg}是不能用的

\section{系统组成及功能说明}\label{sec:sg-1}

次氯酸钠加药装置
由计量泵、脉动阻尼器、背压阀、泄压阀（安全阀）、隔膜压力表、Y型过滤器、反冲洗系统
组成。
  
次氯酸钠加药装置的组成如下图\ref{fig:g02}、图\ref{fig:g03}所示：\\

% 图片模板
      \begin{figure}[h]
         \centering

         \begin{tikzpicture}
            %定义图像
            \node at (0,0) {\includegraphics[height=10cm]{g12.png}};

            %说明
            \node at (5,4) (1) [text_box] {计量泵};
            \draw [arrow] (1) -- (2.2,3);
            %\node at (-5,-9) (2) [text_box] {入药口};
            %\draw [arrow] (2) -- (-1,-7);
            \node at (-4,-4) (3) [text_box] {反冲洗水出口};
            \draw [arrow] (3) -- (-0.9,-3.2);
            \node at (-5,-2) (4) [text_box] {反冲洗出水阀};
            \draw [arrow] (4) -- (-0.9,-2.7);
            \node at (-5,6) (5) [text_box] {隔膜压力表};
            \draw [arrow] (5) -- (-1.4,4);
            \node at (4,-2) (6) [text_box] {加药阀};
            \draw [arrow] (6) -- (0.2,-2.7);
            \node at (-4,0) (6) [text_box] {背压阀};
            \draw [arrow] (6) -- (-0.1,-1.2);
            \node at (-5,2) (7) [text_box] {泄压阀};
            \draw [arrow] (7) -- (-0.6,2.1);
            \node at (4,6) (8) [text_box] {阻尼器};
            \draw [arrow] (8) -- (0.6,4);
            %\node at (5,-7) (9) [text_box] {检修放压排液阀门};
            %\draw [arrow] (9) -- (-3.3,-1.5);
            \node at (4,-4) (10) [text_box] {出药口};
            \draw [arrow] (10) -- (0.2,-3.2);
            \node at (5,0) (11) [text_box] {Y型过滤器};
            \draw [arrow] (11) -- (1.4,-0.6);
            %\node at (6,-5.5) (11) [text_box] {供药阀};
            %\draw [arrow] (11) -- (2.5,-3.8);

			% 辅助线
             \def \xLimit {8};
             \def \yLimit {8};
             %\input{auxiliary_line}

         \end{tikzpicture}
         \caption{次氯酸钠加药装置系统组成-正面}\label{fig:g02}
      \end{figure}




\newpage

\input{picture-system-component-2.tex}

\subsection{加药装置功能特点}
    加药装置有如下功能特点：
    \begin{itemize}
         \item 配备阻尼器及背压阀，消除计量泵投加时产生的脉动。
         \item 配备泄压阀，系统过压时，自动泄压阀避免设备及管道损坏。
         \item 配备隔膜压力表，指示系统当前压力，及显示计量泵工作状态。
         \item 配备反冲洗及检修排液系统，检修前可释放系统压力，排出管道内介质，并使用水等非腐蚀性介质冲洗管道，保证检修人员安全。
    \end{itemize}





\section{设备启动前准备}\label{sec:sg1}
    \subsection{紧固件检查}
    系统启动前，要先对所有的紧固件进行检查。紧固件包括底座固定螺栓、泵固定螺栓、马达连接螺栓以及所有的法兰连接螺栓等。如有松动情况，请将其上紧，如有脱落等情况，请按相应的规格补充相关的螺栓。

    \subsection{管道连接检查}
    检查管道中所有活接接头连接，如有松脱情况需要重新上紧。

    \subsection{检查系统中各阀门状态}
        %\subsubsection{检修排液阀应处于关闭状态}\label{sec:sg2}
        %\input{operation-instruction-drain-valve}
        

        \subsubsection{确认反冲洗控制阀的状态}\label{sec:sg1-1}
        

\begin{figure}[h]
\centering

    \begin{tikzpicture}
        %图片
        \node at (0,3) {\includegraphics[height=10cm]{g12.PNG}};
        \node at (0,-4) {\includegraphics[height=5cm]{g04.PNG}};

        %说明
        \def \itemNumber {1};
        \def \centertPoint {0.6,-3.7};
        \node at (5,-2) (\itemNumber) [text_box] {泵1反冲洗进水阀};
        \draw [arrow] (\itemNumber) -- (\centertPoint);
        \draw (\centertPoint) circle (0.4)[red];
        \def \itemNumber {2};
        \def \centertPoint {-1.4,3.6};
        \node at (-5,2) (\itemNumber) [text_box] {泵1反冲洗出水阀};
        \draw [arrow] (\itemNumber) -- (\centertPoint);
        \draw (\centertPoint) circle (0.4)[red];
        \def \itemNumber {3};
        \def \centertPoint {-0.6,-3.7};
        \node at (-5,-2) (\itemNumber) [text_box] {泵2反冲洗进水阀};
        \draw [arrow] (\itemNumber) -- (\centertPoint);
        \draw (\centertPoint) circle (0.4)[red];
        \def \itemNumber {4};
        \def \centertPoint {1.2,3.6};
        \node at (5,2) (\itemNumber) [text_box] {泵2反冲洗出水阀};
        \draw [arrow] (\itemNumber) -- (\centertPoint);
        \draw (\centertPoint) circle (0.4)[red];
        \def \itemNumber {5};
        \def \centertPoint {-1.1,4.3};
        \node at (-6,4) (\itemNumber) [text_box] {泵1反冲洗隔离阀};
        \draw [arrow] (\itemNumber) -- (\centertPoint);
        \draw (\centertPoint) circle (0.4)[red];
        \def \itemNumber {6};
        \def \centertPoint {0.8,4.3};
        \node at (6,4) (\itemNumber) [text_box] {泵2反冲洗隔离阀};
        \draw [arrow] (\itemNumber) -- (\centertPoint);
        \draw (\centertPoint) circle (0.4)[red];

        % 辅助线
        \def \xLimit {8};
        \def \yLimit {8};
        %\input{auxiliary_line}

    \end{tikzpicture}
\caption{反冲洗控制阀门位置}\label{fig:g02g03}
\end{figure}



反冲洗系统的作用，
是可以使用水等非腐蚀介质冲洗管道内的腐蚀性介质，
在拆卸管道管件时保障维修人员的安全。

\par反冲洗系统 由 反冲洗出水阀、反冲洗进水阀、反冲洗隔离阀组成。
    正常工作时，反冲洗进水阀、反冲洗出水阀应处于关闭状态，反冲洗隔离阀 应处于打开状态。

        

        \subsubsection{其余阀门均处于开启状态}
        除上述特别注明的阀门之外的所有阀门，都应处于开启状态。

\newpage

\section{维护及检修}
%建议联系加药系统经销商及售后服务人员，
%由专业人士进行维护及检修，
%如自行进行维护及检修，
%有可能导致加药系统损坏。
\par维护及检修之前，
请先阅读下文，
以及计量泵的操作说明书。
下文内容为加药装置维护及检修操作顺序。

    \subsection{关闭检修工位的供药阀}\label{sec:sg2}
    \input{operation-instruction-supply-valve.tex}
    
    %\subsection{关闭检修工位的加药阀}\label{sec:sg3}
    %\input{operation-instruction-dosing-valve.tex}

\newpage

    %\subsection{打开相应工位的检修排液阀}
    %\large{\textbf{
    %进行所有维护操作之前，
    %请先打开相应工位的检修排液阀，
    %释放系统内的压力！
    %否则，拆卸相关零部件时，
    %管道内的液体可能会喷射而出，
    %影响维修人员的安全！
    %}}

    \subsection{反冲洗检修工位}
    
加药检修前，
如果有反冲洗条件（反冲洗不影响处理工艺的情况下），
可以使用非腐蚀性（例如水）介质冲洗检修工位计量泵、计量泵附件、阀门、管道，
使操作人员更安全。
\par反冲洗控制阀位置请参考第\pageref{sec:sg1-1}页第\ref{sec:sg1-1}节的图\ref{fig:g02g03}。
以泵1为例，反冲洗操作流程如下：
\begin{enumerate}
    \item 关闭 泵1反冲洗隔离阀。
    \item 打开 泵1反冲洗出水阀。
    \item 打开 泵1反冲洗进水阀，冲洗检修范围内的化学药液。
    \item 关闭 泵1反冲洗进水阀。
    \item 关闭 泵1反冲洗出水阀。
\end{enumerate}



    \subsection{检修设备}
    待检修工位管道内残留化学品排净后，
    可对计量泵、阻尼器、隔膜压力表、背压阀、泄压阀、Y型过液器等设备或元件进行检修操作。
    各设备或元件的检修方法，
    请参阅各设备或元件的使用说明，
    或咨询加药系统经销商或售后服务人员。
    \par对应工位的设备检修完成后，应重新打开 反冲洗隔离阀。
    以泵1工位为例，设备检修完成后，应重新打开 泵1反冲洗隔离阀。

    \subsection{系统恢复运行}
    设备检修完成后，
    请参阅第\ref{sec:sg1}章内容于第\pageref{sec:sg1}页，
    做好设备启动前准备，
    使系统处于待命状态，
    随时恢复正常运行。

\section{反冲洗投加管道}
    
加药装置具有 反冲洗 加药装置 至 投加点之间管道的功能。
\par具体的操作方法如下：
\begin{enumerate}
    \item 根据第\ref{sec:sg2}节于第\pageref{sec:sg2}页的内容，关闭检修工位的供药阀。
    \item 根据第\pageref{fig:g02g03}页图\ref{fig:g02g03}中 反冲洗出水阀的位置，关闭 对应工位的 反冲洗出水阀。
    \item 根据第\pageref{sec:sg-1}页图\ref{fig:g02}中 加药阀位置，保持检修工位的加药阀打开。
    \item 根据第\pageref{fig:g02g03}页图\ref{fig:g02g03}中 反冲洗进水阀的位置，打开 反冲洗进水阀。
    \item 冲洗 加药装置 至 投加点之间管道的管道。
    \item 冲洗 完成后，关闭 反冲洗进水阀。
    \item 打开 检修工位 的 供药阀。
    \item 保持检修工位的加药阀打开，保持检修工位的 反冲洗出水阀 关闭。
\end{enumerate}



\end{document}
